\chapter{Oscilações elétricas livres: RLC}\label{ch:g12:rlc-oscillations}

Gire o botão pesado de um rádio de sótão e as estações desfilam diante do ponteiro: uma
voz, violinos, chiado, outra voz. Atrás do painel não há motor nem computador --- uma
\omterm{def:g11:magnetic-fields:solenoid}{bobina}, um \omterm{def:g12:rc-rl-circuits:capacitor}{capacitor} cujas placas intercaladas o botão gira e um pedaço de fio de
antena. Este capítulo liga os dois tanques de \omterm{def:g10:energy-conservation:energy}{energia} do capítulo anterior
(\cref{ch:g12:rc-rl-circuits}) frente a frente e encontra o pêndulo mais rápido já
construído: carga indo e vindo, até milhões de vezes por segundo, num compasso que o
botão escolhe.

\section{Um pêndulo elétrico: o circuito LC}

Carregue um \omterm{def:g12:rc-rl-circuits:capacitor}{capacitor} $C$ à tensão $U_0$ --- as placas dele guardam então $\pm Q_0$ com
$Q_0 = C U_0$ --- e, em $t = 0$, chaveie-o sobre uma \omterm{def:g11:magnetic-fields:solenoid}{bobina} \emph{ideal} $L$, de
\omterm{def:g11:circuits-and-power:resistance}{resistência} nula. Nenhuma fonte em lugar algum da malha: o que quer que aconteça em
seguida, o circuito faz sozinho.

\begin{omfigure}
\begin{circuitikz}[font=\small]
  \draw (0,0) to[C=$C$, v=$u$] (0,2.4) to[nos, l=$K$] (2.6,2.4)
    -- (3.8,2.4) to[L=$L$, i=$i$] (3.8,0) -- (0,0);
  \node[right, font=\footnotesize] at (0.18,1.5) {$+q$};
  \node[right, font=\footnotesize] at (0.18,0.9) {$-q$};
\end{circuitikz}

{\small O \omterm{prop:g12:rlc-oscillations:equation}{circuito LC}: um \omterm{def:g12:rc-rl-circuits:capacitor}{capacitor} carregado a $Q_0$, uma chave, uma \omterm{def:g11:magnetic-fields:solenoid}{bobina} ideal.
Fechar $K$ deixa a carga escorrer pela \omterm{def:g11:magnetic-fields:solenoid}{bobina} --- e a \omterm{def:g11:magnetic-fields:solenoid}{bobina}, que detesta mudança, não
a deixa parar.}
\end{omfigure}

\begin{proposition}[A equação do LC]\label{prop:g12:rlc-oscillations:equation}
\index{circuito LC}
Na malha LC ideal a carga $q(t)$ da placa de cima obedece a
\[
L\,\frac{d^2q}{dt^2} + \frac{q}{C} = 0,
\qquad\text{ou seja,}\qquad
\frac{d^2q}{dt^2} = -\frac{1}{LC}\,q .
\]
\end{proposition}

\begin{proof}
Ao longo da malha as tensões somam zero (\cref{prop:g11:circuits-and-power:laws}):
$u_L + u_C = 0$. O \omterm{def:g12:rc-rl-circuits:capacitor}{capacitor} dá $u_C = q/C$ e a \omterm{def:g11:magnetic-fields:solenoid}{bobina}, $u_L = L\,di/dt$
(\cref{def:g12:rc-rl-circuits:inductor}), com $i = dq/dt$ contando a carga que chega à
placa de cima (\cref{prop:g12:rc-rl-circuits:cdudt}); substitua.
\end{proof}

\begin{theorem}[Oscilações livres do circuito LC]\label{thm:g12:rlc-oscillations:oscillation}
\index{oscilações elétricas}
Solto em $t = 0$ com $q(0) = Q_0$ e $i(0) = 0$, o circuito oscila:
\[
q(t) = Q_0 \cos\!\left(\frac{2\pi t}{T_0}\right),
\qquad
i(t) = -I_0 \sin\!\left(\frac{2\pi t}{T_0}\right),
\qquad
T_0 = 2\pi\sqrt{LC},
\]
com corrente de pico $I_0 = 2\pi Q_0/T_0 = Q_0/\sqrt{LC}$.
\end{theorem}

\begin{proof}
Verifique por substituição. Derivando,
$i = dq/dt = -\frac{2\pi}{T_0} Q_0 \sin(2\pi t/T_0)$ --- o $i(t)$ enunciado --- e de
novo $\frac{d^2q}{dt^2} = -(\frac{2\pi}{T_0})^2 q$, que é igual a $-q/(LC)$ exatamente
quando $(2\pi/T_0)^2 = 1/(LC)$, ou seja, $T_0 = 2\pi\sqrt{LC}$. E $q(0) = Q_0$,
$i(0) = 0$. Que nenhuma \emph{outra} função se ajusta à equação e à partida fica
admitido, exatamente como no caso da mola
(\cref{rem:g12:oscillators-and-time:uniqueness}).
\end{proof}

\begin{definition}[Período e frequência próprios]\label{def:g12:rlc-oscillations:period}
O \emph{período próprio}\index{período próprio!de um circuito LC} de um \omterm{prop:g12:rlc-oscillations:equation}{circuito LC} é
$T_0 = 2\pi\sqrt{LC}$; a \emph{frequência própria}\index{frequência própria!de um
circuito LC} dele é
\[
f_0 = \frac{1}{T_0} = \frac{1}{2\pi\sqrt{LC}} .
\]
Conferência dimensional: $\qty{1}{H} = \qty{1}{V.s/A}$ e $\qty{1}{F} = \qty{1}{C/V}$,
de modo que $LC$ carrega $\unit{s} \times \unit{C/A} = \unit{s^2}$ --- $\sqrt{LC}$ é um
tempo, como tem de ser.
\end{definition}

\begin{example}[Primeiros números]\label{ex:g12:rlc-oscillations:firstnumbers}
$L = \qty{10}{mH}$ com $C = \qty{1.0}{\micro F}$: $\sqrt{LC} = \qty{1.0e-4}{s}$, logo
$T_0 = \qty{0.63}{ms}$ e $f_0 \approx \qty{1.6}{kHz}$ --- um tom audível.
$L = \qty{1.0}{mH}$ com $C = \qty{100}{pF}$: $T_0 = \qty{2.0}{\micro s}$,
$f_0 \approx \qty{0.50}{MHz}$ --- uma radiofrequência. Só o \emph{produto} $LC$ fixa o
compasso.
\end{example}

\begin{remark}[Um quarto de período fora de compasso]\label{rem:g12:rlc-oscillations:quarter}
A corrente está um quarto de \omterm{def:g12:oscillators-and-time:oscillator}{período} defasada em relação à carga: $i = 0$ quando
$q = \pm Q_0$ (placas cheias, escoamento se invertendo) e $i = \pm I_0$ quando $q = 0$
(placas vazias, escoamento a todo vapor) --- exatamente o pêndulo, imóvel nos extremos e
mais rápido embaixo.
\end{remark}

\begin{omfigure}
\begin{tikzpicture}
\begin{axis}[omaxis, width=9.2cm, height=5.0cm, xlabel={$t$},
    xmin=0, xmax=2, ymin=-1.4, ymax=1.75, xtick={0,0.5,1,1.5,2},
    xticklabels={$0$,$T_0/2$,$T_0$,$3T_0/2$,$2T_0$},
    ytick={-1,0,1}, legend pos=north east]
  \addplot[omDef, thick, domain=0:2, samples=200] {cos(360*x)};
  \addplot[omThm, thick, domain=0:2, samples=200] {-sin(360*x)};
  \legend{$q/Q_0$, $i/I_0$}
\end{axis}
\end{tikzpicture}

{\small Carga e corrente do \omterm{prop:g12:rlc-oscillations:equation}{circuito LC} livre: a corrente é máxima toda vez que a carga
cruza o zero, um quarto de \omterm{def:g12:oscillators-and-time:oscillator}{período} fora de compasso.}
\end{omfigure}

\section{A valsa das energias}

\begin{proposition}[Conservação da energia armazenada]\label{prop:g12:rlc-oscillations:energy}
\index{energia!de um circuito LC}
No \omterm{prop:g12:rlc-oscillations:equation}{circuito LC} ideal a \omterm{def:g10:energy-conservation:energy}{energia} total armazenada
(\cref{prop:g12:rc-rl-circuits:energy})
\[
E \;=\; E_C + E_L \;=\; \frac{q^2}{2C} + \tfrac12\,L i^2
\;=\; \frac{Q_0^2}{2C} \;=\; \tfrac12\,L I_0^2
\]
é constante: inteiramente no \omterm{def:g12:rc-rl-circuits:capacitor}{capacitor} quando $q = \pm Q_0$, inteiramente na \omterm{def:g11:magnetic-fields:solenoid}{bobina}
quando $i = \pm I_0$.
\end{proposition}

\begin{proof}
Derive: $\frac{dE}{dt} = \frac{q}{C}\frac{dq}{dt} + Li\frac{di}{dt} =
i\left(\frac{q}{C} + L\frac{d^2q}{dt^2}\right) = 0$ pela
\cref{prop:g12:rlc-oscillations:equation}. O valor dela se lê em $t = 0$ (tudo no
\omterm{def:g12:rc-rl-circuits:capacitor}{capacitor}) e de novo no instante em que tudo é corrente, $q = 0$ --- de onde
$\tfrac12 L I_0^2 = Q_0^2/2C$, recuperando $I_0 = Q_0/\sqrt{LC}$ só com \omterm{def:g10:energy-conservation:energy}{energia}.
\end{proof}

\begin{omfigure}
\begin{tikzpicture}
\begin{axis}[omaxis, width=8.6cm, height=4.6cm, xlabel={$t/T_0$},
    ylabel={energia $/\,E$}, xmin=0, xmax=1, ymin=0, ymax=1.3,
    legend pos=south east]
  \addplot[omDef, thick, domain=0:1, samples=120] {cos(360*x)^2};
  \addplot[omProp, thick, domain=0:1, samples=120] {sin(360*x)^2};
  \addplot[omThm, dashed, thick, domain=0:1, samples=2] {1};
  \legend{$E_C$, $E_L$, $E$}
\end{axis}
\end{tikzpicture}

{\small Um \omterm{def:g12:oscillators-and-time:oscillator}{período} da valsa: a \omterm{def:g10:energy-conservation:energy}{energia} vai e volta das placas para a \omterm{def:g11:magnetic-fields:solenoid}{bobina} duas vezes
por \omterm{def:g12:oscillators-and-time:oscillator}{período}, com a soma cravada em $Q_0^2/2C$ --- a mesmíssima figura do \omterm{prop:g12:oscillators-and-time:spring-period}{oscilador
massa--mola}, com outros rótulos.}
\end{omfigure}

\begin{example}[A valsa, medida]\label{ex:g12:rlc-oscillations:waltz}
$C = \qty{1.0}{\micro F}$ carregado a $U_0 = \qty{10}{V}$, chaveado sobre
$L = \qty{10}{mH}$: \omterm{def:g10:energy-conservation:energy}{energia} $E = \tfrac12 C U_0^2 = \qty{5.0e-5}{J}$, \omterm{def:g12:oscillators-and-time:oscillator}{período}
$T_0 = \qty{0.63}{ms}$ e uma corrente de pico $I_0 = U_0\sqrt{C/L} = \qty{0.10}{A}$
atingida um quarto de \omterm{def:g12:oscillators-and-time:oscillator}{período} (\qty{0.16}{ms}) depois que a chave fecha.
\end{example}

\section{Circuitos reais: oscilações amortecidas}

Toda \omterm{def:g11:magnetic-fields:solenoid}{bobina} real é enrolada com fio resistivo: a malha honesta é um circuito
\emph{RLC} série, e a \omterm{def:g11:circuits-and-power:resistance}{resistência} taxa a valsa.

\begin{definition}[Regimes pseudoperiódico e aperiódico]\label{def:g12:rlc-oscillations:regimes}
\index{circuito RLC}\index{regime pseudoperiódico}\index{regime aperiódico}
Num circuito RLC série o resistor dissipa \omterm{def:g10:energy-conservation:energy}{energia} por \omterm{def:g11:circuits-and-power:joule}{efeito Joule}
(\cref{ch:g11:circuits-and-power}): as \omterm{def:g12:oscillators-and-time:oscillator}{oscilações livres} são \emph{\omterm{def:g12:oscillators-and-time:damping}{amortecidas}}. Dois
regimes, lidos num \omterm{def:g10:signals-and-waves:oscilloscope}{osciloscópio}:
\begin{itemize}
  \item \emph{pseudoperiódico} ($R$ pequeno): o circuito ainda oscila dentro de um
        envelope que encolhe, com um tempo de repetição --- o
        \emph{pseudoperíodo}\index{pseudoperíodo!elétrico} --- próximo de
        $T_0 = 2\pi\sqrt{LC}$ quando o \omterm{def:g12:oscillators-and-time:damping}{amortecimento} é fraco;
  \item \emph{aperiódico} ($R$ grande): a carga volta rastejando a zero sem nunca
        passar do ponto --- oscilação nenhuma.
\end{itemize}
O espelho exato do \omterm{def:g12:oscillators-and-time:oscillator}{oscilador mecânico} amortecido
(\cref{def:g12:oscillators-and-time:damping}).
\end{definition}

\begin{omfigure}
\begin{tikzpicture}
\begin{axis}[omaxis, width=9.2cm, height=5.2cm,
    xlabel={$t$ (\unit{\micro s})}, ylabel={$u$ (\unit{V})},
    xmin=0, xmax=10, ymin=-6.6, ymax=8.2, legend pos=north east]
  \addplot[omDef, thick, domain=0:10, samples=300]
    {6*exp(-x/3)*cos(deg(2*pi*x/1.6))};
  \addplot[omProp, thick, domain=0:10, samples=100] {6*exp(-x/1.2)};
  \legend{pseudoperiódico ($R$ pequeno), aperiódico ($R$ grande)}
  \addplot[omThm, dashed, domain=0:10, samples=80, forget plot]
    {6*exp(-x/3)};
  \addplot[omThm, dashed, domain=0:10, samples=80, forget plot]
    {-6*exp(-x/3)};
\end{axis}
\end{tikzpicture}

{\small Os dois destinos de um \omterm{def:g12:rlc-oscillations:regimes}{circuito RLC} real: oscilações que morrem dentro de um
envelope exponencial (tracejado), ou um rastejo até zero sem oscilação alguma. Com
\omterm{def:g12:oscillators-and-time:damping}{amortecimento} fraco, o \omterm{def:g12:rlc-oscillations:regimes}{pseudoperíodo} continua $\approx T_0$.}
\end{omfigure}

\begin{remark}[Manter as oscilações]\label{rem:g12:rlc-oscillations:maintain}
\index{oscilações mantidas}
Para manter a \omterm{def:g12:oscillators-and-time:oscillator}{amplitude} viva, injete no circuito exatamente a \omterm{def:g10:energy-conservation:energy}{energia} que a \omterm{def:g11:circuits-and-power:resistance}{resistência}
queima, uma vez por ciclo e \emph{em compasso} com o balanço: um amplificador faz pela
malha RLC o que o empurrão do escapamento faz pelo pêndulo
(\cref{ch:g12:oscillators-and-time}). Todo relógio de quartzo e todo transmissor de
rádio é um \omterm{def:g12:oscillators-and-time:oscillator}{oscilador} mantido desse tipo, rodando na \omterm{def:g12:rlc-oscillations:period}{frequência própria} $f_0$ dele.
\end{remark}

\section{O gêmeo mecânico}

\begin{proposition}[O dicionário elétrico--mecânico]\label{prop:g12:rlc-oscillations:analogy}
\index{analogia elétrica--mecânica}
A equação do LC é a equação massa--mola
(\cref{prop:g12:oscillators-and-time:spring-period}) palavra por palavra:
\begin{center}\small
\begin{tabular}{lcl}
\omterm{prop:g12:rlc-oscillations:equation}{circuito LC} & & carrinho massa--mola\\[2pt]
carga $q$ & $\longleftrightarrow$ & posição $x$\\
corrente $i = dq/dt$ & $\longleftrightarrow$ & velocidade $v = dx/dt$\\
\omterm{def:g12:rc-rl-circuits:inductor}{indutância} $L$ & $\longleftrightarrow$ & massa $m$\\
inverso da \omterm{def:g12:rc-rl-circuits:capacitor}{capacitância} $1/C$ & $\longleftrightarrow$ & \omterm{def:g12:oscillators-and-time:spring}{constante elástica} $k$\\
\omterm{def:g11:circuits-and-power:resistance}{resistência} $R$ & $\longleftrightarrow$ & \omterm{def:g10:inertia:inventory}{atrito}\\
$L\,\frac{d^2q}{dt^2} + \frac{q}{C} = 0$ & $\longleftrightarrow$ &
  $m\,\frac{d^2x}{dt^2} + kx = 0$\\[2pt]
$T_0 = 2\pi\sqrt{LC}$ & $\longleftrightarrow$ &
  $T_0 = 2\pi\sqrt{m/k}$\\[2pt]
$E_C = q^2/2C$ & $\longleftrightarrow$ & $E_p = \tfrac12 kx^2$\\
$E_L = \tfrac12 Li^2$ & $\longleftrightarrow$ & $E_k = \tfrac12 mv^2$
\end{tabular}
\end{center}
Mesma equação, mesmas soluções: todo resultado sobre um dos \omterm{def:g12:oscillators-and-time:oscillator}{osciladores} se traduz na
hora para a outra língua.
\end{proposition}

\begin{proof}
Compare a \cref{prop:g12:rlc-oscillations:equation} com $m\,d^2x/dt^2 = -kx$:
substituir $q \to x$, $L \to m$ e $1/C \to k$ transforma uma na outra, e
$LC \to m/k$ transforma $2\pi\sqrt{LC}$ em $2\pi\sqrt{m/k}$.
\end{proof}

\begin{example}[Sintonizar um rádio]\label{ex:g12:rlc-oscillations:tuning}
A antena joga um sussurro de \emph{todas} as estações dentro de uma malha LC, mas só a
que transmite perto de $f_0$ empurra em compasso e vai crescendo --- \omterm{def:g12:oscillators-and-time:resonance}{ressonância}
(\cref{def:g12:oscillators-and-time:resonance}). O botão gira as placas de um \omterm{def:g12:rc-rl-circuits:capacitor}{capacitor}
variável: mudar $C$ desloca $f_0 = 1/(2\pi\sqrt{LC})$, e o ponteiro entrega ao fone uma
estação de cada vez. Do lado da transmissão, um \omterm{def:g12:oscillators-and-time:oscillator}{oscilador} mantido
(\cref{rem:g12:rlc-oscillations:maintain}) fixa a \omterm{def:g12:oscillators-and-time:oscillator}{frequência} de transmissão; num
relógio, um cristal de quartzo --- um gêmeo mecânico de rigidez de tirar o fôlego ---
faz o mesmo papel.
\end{example}

\section{Exercícios}

\begin{exercise}[$\star$]\label{exo:g12:rlc-oscillations:1}
Calcule $T_0$ e $f_0$ para (a) $L = \qty{10}{mH}$, $C = \qty{1.0}{\micro F}$; (b)
$L = \qty{1.0}{mH}$, $C = \qty{100}{pF}$. Por que fator diferem as \omterm{def:g12:oscillators-and-time:oscillator}{frequências}?
\end{exercise}

\begin{exercise}[$\star$]\label{exo:g12:rlc-oscillations:2}
Usando $\qty{1}{H} = \qty{1}{V.s/A}$ e $\qty{1}{F} = \qty{1}{C/V}$, mostre, \omterm{def:g10:orders-of-magnitude:unit}{unidade} por
\omterm{def:g10:orders-of-magnitude:unit}{unidade}, que $\sqrt{LC}$ é um tempo. Por que o $2\pi$ não muda nada?
\end{exercise}

\begin{exercise}[$\star$]\label{exo:g12:rlc-oscillations:3}
Um \omterm{def:g12:rc-rl-circuits:capacitor}{capacitor} de \qty{470}{nF} é carregado a \qty{6.0}{V}. Calcule $Q_0$ e a \omterm{def:g10:energy-conservation:energy}{energia}
armazenada; descreva, sem calcular, o que acontece quando ele é chaveado sobre uma
\omterm{def:g11:magnetic-fields:solenoid}{bobina} ideal.
\end{exercise}

\begin{exercise}[$\star$]\label{exo:g12:rlc-oscillations:4}
Um \omterm{prop:g12:rlc-oscillations:equation}{circuito LC} segue $q(t) = 2.0\cos(2\pi t/T_0)$ em microcoulombs, com
$T_0 = \qty{4.0}{ms}$. Leia $Q_0$; calcule $f_0$, $q(T_0/2)$, $i(0)$ e a corrente de
pico $I_0 = 2\pi Q_0/T_0$.
\end{exercise}

\begin{exercise}[$\star$]\label{exo:g12:rlc-oscillations:5}
Em que instantes do ciclo a \omterm{def:g10:energy-conservation:energy}{energia} está inteiramente no \omterm{def:g12:rc-rl-circuits:capacitor}{capacitor}? E inteiramente na
\omterm{def:g11:magnetic-fields:solenoid}{bobina}? Qual é a contrapartida mecânica de cada momento para um pêndulo?
\end{exercise}

\begin{exercise}[$\star\star$]\label{exo:g12:rlc-oscillations:6}
Verifique por substituição que $q(t) = Q_0\cos(2\pi t/T_0)$ satisfaz
$L\,d^2q/dt^2 + q/C = 0$ exatamente quando $T_0 = 2\pi\sqrt{LC}$.
\end{exercise}

\begin{exercise}[$\star\star$]\label{exo:g12:rlc-oscillations:7}
Uma estação de \omterm{def:g10:signals-and-waves:wave}{ondas} longas transmite a \qty{198}{kHz}. Com $C = \qty{220}{pF}$, que
\omterm{def:g12:rc-rl-circuits:inductor}{indutância} sintoniza nela?
\end{exercise}

\begin{exercise}[$\star\star$]\label{exo:g12:rlc-oscillations:8}
$C = \qty{1.0}{\micro F}$ carregado a \qty{10}{V} se descarrega em $L = \qty{10}{mH}$.
Calcule a \omterm{def:g10:energy-conservation:energy}{energia} armazenada, a corrente de pico (pela \omterm{thm:g10:energy-conservation:conservation}{conservação da energia}) e o
instante em que ela é atingida pela primeira vez.
\end{exercise}

\begin{exercise}[$\star\star$]\label{exo:g12:rlc-oscillations:9}
Explique por que $E_C$ e $E_L$ oscilam na \omterm{def:g12:oscillators-and-time:oscillator}{frequência} $2f_0$, e não $f_0$. Mostre que em
$t = T_0/8$ (partindo de $q = Q_0$) a \omterm{def:g10:energy-conservation:energy}{energia} está dividida exatamente ao meio.
\end{exercise}

\begin{exercise}[$\star\star$]\label{exo:g12:rlc-oscillations:10}
Num \omterm{def:g10:signals-and-waves:oscilloscope}{oscilograma} de um \omterm{def:g12:rlc-oscillations:regimes}{circuito RLC} fracamente amortecido, máximos sucessivos marcam
\qty{4.0}{V} e depois \qty{3.0}{V}. Supondo que a razão permaneça constante, preveja o
sexto máximo depois do de \qty{4.0}{V}. Que fração da \emph{\omterm{def:g10:energy-conservation:energy}{energia}} sobrevive a cada
\omterm{def:g12:rlc-oscillations:regimes}{pseudoperíodo}? E o que se pode dizer do \omterm{def:g12:rlc-oscillations:regimes}{pseudoperíodo} em si?
\end{exercise}

\begin{exercise}[$\star\star$]\label{exo:g12:rlc-oscillations:11}
Um \omterm{prop:g12:rlc-oscillations:equation}{circuito LC} tem $L = \qty{0.50}{H}$ e $C = \qty{20}{\micro F}$. Calcule $T_0$; e
depois, usando o dicionário (\cref{prop:g12:rlc-oscillations:analogy}), encontre a
\omterm{def:g12:oscillators-and-time:spring}{constante elástica} $k$ que dá a um carrinho de \qty{0.50}{kg} o mesmo \omterm{def:g12:oscillators-and-time:oscillator}{período}, e
confira.
\end{exercise}

\begin{exercise}[$\star\star\star$]\label{exo:g12:rlc-oscillations:12}
Um sintonizador usa $L = \qty{0.20}{mH}$ e um \omterm{def:g12:rc-rl-circuits:capacitor}{capacitor} variável que varre de
\qtyrange{50}{500}{pF}. Calcule as duas \omterm{def:g12:oscillators-and-time:oscillator}{frequências} extremas. Mostre que
$f_0 \propto 1/\sqrt{C}$, de modo que dez vezes em $C$ dá apenas
$\sqrt{10} \approx 3.2$ em \omterm{def:g12:oscillators-and-time:oscillator}{frequência}.
\end{exercise}

\begin{exercise}[$\star\star\star$]\label{exo:g12:rlc-oscillations:13}
Uma \omterm{def:g11:magnetic-fields:solenoid}{bobina} real faz o circuito perder $5.0\%$ da \omterm{def:g10:energy-conservation:energy}{energia} dele a cada \omterm{def:g12:rlc-oscillations:regimes}{pseudoperíodo}.
Depois de quantos \omterm{def:g12:oscillators-and-time:oscillator}{períodos} a \omterm{def:g10:energy-conservation:energy}{energia} caiu à metade? Qual é a \omterm{def:g12:oscillators-and-time:oscillator}{amplitude} então, como
fração da inicial? A $f_0 = \qty{1.0}{MHz}$, quanto tempo é isso em microssegundos?
\end{exercise}

\begin{exercise}[$\star\star\star$]\label{exo:g12:rlc-oscillations:14}
O \omterm{def:g12:oscillators-and-time:oscillator}{oscilador} de um relógio armazena $E = \qty{1.0e-9}{J}$ e perde $2.0\%$ disso por
ciclo, a $f_0 = \qty{32768}{Hz}$. Que \omterm{def:g11:work-of-force:power}{potência} média o circuito de manutenção precisa
fornecer? Por que o empurrão tem de chegar em compasso com a oscilação?
\end{exercise}

\begin{exercise}[$\star\star\star$]\label{exo:g12:rlc-oscillations:15}
Depois de cada borda de um \omterm{def:g10:signals-and-waves:signal}{sinal} quadrado, a fiação de um experimentador ``repica'' a
\qty{25}{MHz}; a \omterm{def:g12:rc-rl-circuits:inductor}{indutância} parasita é de cerca de \qty{1.0}{\micro H}. Estime a
\omterm{def:g12:rc-rl-circuits:capacitor}{capacitância} parasita responsável. Para que regime um pequeno resistor em série deveria
empurrar o circuito, e por que isso cura o repique?
\end{exercise}

\section{Problema: montar um rádio de galena}

\begin{problem}[{Problema de fim de semana --- montar um rádio de galena: uma bobina, um
capacitor variável, um diodo e um fone --- sem pilha, sem amplificador, e o jornal da
noite chega com a energia da própria onda}]
\label{pb:g12:rlc-oscillations:1}
Um kit do sótão dos avós: uma \omterm{def:g11:magnetic-fields:solenoid}{bobina} $L = \qty{0.20}{mH}$ enrolada num tubo de papelão,
um \omterm{def:g12:rc-rl-circuits:capacitor}{capacitor} variável cujas placas giram de \qty{20}{pF} a \qty{480}{pF}, um diodo de
cristal, um fone de ouvido e \qty{20}{m} de fio de antena. As estações a pegar
transmitem entre \qty{530}{kHz} e \qty{1700}{kHz}.

\medskip
\noindent\textbf{Parte I --- A malha de sintonia.}
\begin{enumerate}
  \item O \omterm{def:g12:rc-rl-circuits:capacitor}{capacitor}, carregado, é deixado descarregar na \omterm{def:g11:magnetic-fields:solenoid}{bobina}: escreva a lei das
        malhas e transforme-a na equação diferencial de $q(t)$.
  \item Verifique por substituição que $q(t) = Q_0\cos(2\pi t/T_0)$ a resolve, desde
        que $T_0 = 2\pi\sqrt{LC}$.
  \item Confira por análise dimensional que $\sqrt{LC}$ é um tempo.
  \item Botão ajustado em $C = \qty{330}{pF}$: calcule $T_0$ e $f_0$. A malha está
        sintonizada dentro da faixa?
  \item Encontre a \omterm{def:g12:rc-rl-circuits:capacitor}{capacitância} que sintoniza exatamente \qty{1.00}{MHz}.
\end{enumerate}

\medskip
\noindent\textbf{Parte II --- O \omterm{def:g12:rc-rl-circuits:capacitor}{capacitor} variável.}
\begin{enumerate}[resume]
  \item Calcule a \omterm{def:g12:rc-rl-circuits:capacitor}{capacitância} necessária para \qty{530}{kHz}.
  \item Calcule a \omterm{def:g12:rc-rl-circuits:capacitor}{capacitância} necessária para \qty{1700}{kHz}.
  \item Calcule as duas \omterm{def:g12:oscillators-and-time:oscillator}{frequências} que o \omterm{def:g12:rc-rl-circuits:capacitor}{capacitor} de \qtyrange{20}{480}{pF} do kit
        de fato alcança. Ele cobre a faixa inteira?
  \item Mostre que $f_0 \propto 1/\sqrt{C}$ com $L$ fixo; por que fator $C$ precisa
        mudar para dobrar $f_0$?
  \item As placas se abrem linearmente, de modo que $C$ cai de forma uniforme conforme
        o botão gira. Explique por que as estações então se amontoam numa das pontas do
        dial --- e em qual delas.
\end{enumerate}

\medskip
\noindent\textbf{Parte III --- \omterm{def:g10:energy-conservation:energy}{Energia} e \omterm{def:g12:oscillators-and-time:damping}{amortecimento}.}
Sintonizada em \qty{1.00}{MHz} ($C = \qty{127}{pF}$), a antena carrega por um instante
o \omterm{def:g12:rc-rl-circuits:capacitor}{capacitor} a $U_0 = \qty{20}{mV}$.
\begin{enumerate}[resume]
  \item Calcule a carga $Q_0$ e a \omterm{def:g10:energy-conservation:energy}{energia} armazenada $E$.
  \item Calcule a corrente de pico $I_0$, usando a \omterm{thm:g10:energy-conservation:conservation}{conservação da energia}.
  \item A \omterm{def:g11:circuits-and-power:resistance}{resistência} da \omterm{def:g11:magnetic-fields:solenoid}{bobina} dissipa $4.0\%$ da \omterm{def:g10:energy-conservation:energy}{energia} armazenada por ciclo: depois
        de quantos ciclos a \omterm{def:g10:energy-conservation:energy}{energia} caiu à metade?
  \item Quanto tempo é isso em microssegundos? Por que esse ``repique'' sozinho não
        consegue tocar o jornal da noite?
  \item Explique por que a \omterm{def:g10:signals-and-waves:wave}{onda} que chega mantém a malha balançando apenas quando ela
        está sintonizada naquela estação --- e o que mantém a oscilação no
        \emph{transmissor} (\cref{rem:g12:rlc-oscillations:maintain}).
\end{enumerate}

\medskip
\noindent\textbf{Parte IV --- O gêmeo mecânico.}
\begin{enumerate}[resume]
  \item Usando o dicionário (\cref{prop:g12:rlc-oscillations:analogy}) com um carrinho
        de massa $m = \qty{0.20}{kg}$, calcule a \omterm{def:g12:oscillators-and-time:spring}{constante elástica} $k$ da mola que
        corresponde ao circuito sintonizado.
  \item Uma mola dura de carro tem $k \approx \qty{5e4}{N/m}$. Comente: um sistema
        massa--mola de bancada consegue oscilar em megahertz?
  \item Fique, em vez disso, com um $k = \qty{100}{N/m}$ razoável: que massa oscila a
        \qty{1.00}{MHz}?
  \item Que objeto real é um \omterm{def:g12:oscillators-and-time:oscillator}{oscilador mecânico} assim, leve como pluma e ultrarrígido, e
        onde você o encontrou antes, este ano
        (\cref{ch:g12:oscillators-and-time})?
  \item Ficha técnica, em uma frase: a faixa que o seu rádio varre, a \omterm{def:g12:rc-rl-circuits:capacitor}{capacitância} que
        pega \qty{1.00}{MHz}, quanto dura um repique solitário e o que toca o jornal
        assim mesmo.
\end{enumerate}
\end{problem}
